\section*{NS-96: divisor squares before and after the old-space refit}

\paragraph{Scope.}
A genuine pointwise divisor square uses least common multiples, not the
Dirichlet convolution square of NS86. We test finite positive semidefinite
mixtures of such squares. As complete coefficient rules, either global sign
has a uniform positive approximation-error floor. After an unrestricted
old-space refit, their projected cone is exactly a signed linear space.
Thus this positivity condition alone supplies no sign for the actual
residual numerator. None of these statements excludes convergence of the
unrestricted integer approximation, or supplies its missing lower gain.
The manuscript remains v1.66.

\paragraph{Prior art and normalization.}
The divisor-square/LCM algebra and its diagonalization are classical
Selberg-sieve tools; see Steve Lester, \emph{The Selberg sieve: lecture 2}
(May 7, 2015), pp. 1 and 3, hosted at
\url{https://www.math.tau.ac.il/~rudnick/courses/sieves2015/SelbergSieve_lecture_2.pdf}.
We claim no novelty for those identities. All arguments needed here are
given below. Keep the fixed NS61/78/86 normalization
\[
 \mathcal H=L^2((0,\infty),dt/t^2),\quad
 A(z)=\int_0^z\{u\}\,du/u,\quad a_n(t)=A(t/n),\quad
 \tau(t)=\log_+t.
\]
For any finite real coefficient sequence, set
$r=\tau-\sum c_n a_n$, $E=\|r\|^2$ and $\eta=\sum c_n/n$.
Then $r(t)=-\eta t$ below 1. No optimality hypothesis is imposed unless
stated. In particular, a lower bound for a restricted coefficient rule is
not a lower bound for the optimized error $E_N$.

\subsection*{1. The actual divisor square}
Fix $N\geq1$ and a real symmetric matrix $\Gamma\succeq0$ of size $N$.
Write $[i,j]=\operatorname{lcm}(i,j)$ and group coefficients by
\[
 q_\ell=\sum_{\substack{i,j\leq N\\[i,j]=\ell}}\Gamma_{ij},
 \qquad S_N=\{[i,j]:i,j\leq N\}\subseteq[1,N^2].
 \tag{96.1}
\]
For every positive integer $m$, not merely $m\leq N^2$,
\[
 \sum_{\ell\mid m}q_\ell=v_m^T\Gamma v_m\geq0,
 \qquad (v_m)_i=\boldsymbol1_{i\mid m}.
 \tag{96.2}
\]
Indeed, $[i,j]\mid m$ if and only if both $i\mid m$ and $j\mid m$.
For rank one $\Gamma=\lambda\lambda^T$ this is
$(\sum_{i\mid m,\ i\leq N}\lambda_i)^2$.
The individual $q_\ell$ can have either sign. This is neither a cone of
nonnegative atom coefficients nor NS64's positive adjacent-difference cone.

For every real $s>0$, define $Q(s)=\sum_\ell q_\ell\ell^{-s}$.
The reciprocal LCM matrix $K_s=([i,j]^{-s})$ is positive definite. To see
this without a numerical eigenvalue calculation, put
\[
 J_s(d)=d^s\prod_{p\mid d}(1-p^{-s})>0,
 \quad \sum_{d\mid k}J_s(d)=k^s.
\]
The product identity follows prime-power by prime-power. Hence
\[
 \sum_{i,j\leq N}\lambda_i\lambda_j[i,j]^{-s}
 =\sum_{d\leq N}J_s(d)
     \left(\sum_{\substack{i\leq N\\d\mid i}}\lambda_i i^{-s}\right)^2.
 \tag{96.3}
\]
The divisor-incidence matrix is triangular with nonzero diagonal, so the
right side is strictly positive for nonzero $\lambda$. Decomposing
$\Gamma$ into rank-one positive terms proves
\[
 Q(s)=\operatorname{tr}(\Gamma K_s)\geq0.
 \tag{96.4}
\]
This matrix is not the physical arithmetic Gram matrix.

\subsection*{2. A complete real Mellin separator}
For $1/2<s<1$, let $C(s)=\sum c_n n^{-s}$. Complete integration gives
\[
 \mathcal L_s(r):=\int_0^\infty r(t)t^{-s-1}\,dt
       =\frac{1+\zeta(s)C(s)}{s^2}.
 \tag{96.5}
\]
Here $\int\tau(t)t^{-s-1}dt=1/s^2$. Also
\[
 \int_0^\infty\{t\}t^{-s-1}dt=-\zeta(s)/s,
 \qquad
 \int_0^\infty A(t/n)t^{-s-1}dt=-\zeta(s)n^{-s}/s^2.
\]
For completeness, the first formula follows for $\Re s>1$ by subtracting
$\int_1^\infty\lfloor t\rfloor t^{-s-1}dt=\zeta(s)/s$ from
$\int_1^\infty t^{-s}dt$, then continuing the convergent fractional-part
integral on $[1,\infty)$ to $\Re s>0$ and adding $\int_0^1 t^{-s}dt$.
The second follows by nonnegative Fubini for real $0<s<1$ and a dilation.
This also proves $\zeta(s)<0$ on $(0,1)$. Endpoint integrability uses
$r(t)=O(t)$ at zero and $r(t)=O(1+\log t)$ at infinity.

The bare Mellin functional is not bounded on all of $\mathcal H$. We use
only the subspace of functions linear in $t$ on $(0,1)$, which contains
every residual here. On that subspace its bounded representative is
\[
 k_s(t)=\begin{cases}t/(1-s),&0<t<1,\\t^{1-s},&t\geq1,\end{cases}
 \quad \mathcal L_s(r)=\langle r,k_s\rangle,
 \quad \|k_s\|^2=\frac1{(1-s)^2}+\frac1{2s-1}
       =\frac{s^2}{(1-s)^2(2s-1)}.
 \tag{96.6}
\]
In particular no exterior term has been dropped.

\paragraph{Negative complete coefficient rule.}
If the \emph{entire} coefficient sequence is $c_\ell=-q_\ell$, (96.4)--(96.6)
and Cauchy--Schwarz imply, for every $N$ and every $\Gamma\succeq0$,
\[
 E\geq\frac{(1-s)^2(2s-1)}{s^6};
 \qquad s=4/7\quad\Longrightarrow\quad
 \boxed{E\geq3087/4096}.
 \tag{96.7}
\]
Indeed $\mathcal L_s(r)=(1-\zeta(s)Q(s))/s^2\geq1/s^2$.
The bound is exact and independent of a numerical cutoff or of $N$.

\paragraph{Positive complete coefficient rule.}
If instead $c_\ell=q_\ell$, then $c_1=\Gamma_{11}\geq0$.
On $(0,2)$ one has exactly
$r(t)=-\eta t+(1+c_1)\log_+t$.
Restricting the nonnegative norm integral to that interval and minimizing
over arbitrary $\eta\in\mathbb R$ gives
\[
 E\geq(1+c_1)^2\kappa_0\geq\kappa_0>1/32,\qquad
 \kappa_0=1-\log2-\tfrac12(\log2)^2-\tfrac18(\log2)^4.
 \tag{96.8}
\]
In fact the local squared norm is
$2\eta^2-\eta(1+c_1)(\log2)^2+
(1+c_1)^2(1-\log2-(\log2)^2/2)$.
To check positivity using rationals alone, the series
$\log2=2\sum_{k\geq0}3^{-2k-1}/(2k+1)$ gives
$\log2<2/3+1/36=25/36$. The function
$f(x)=1-x-x^2/2-x^4/8$ is strictly decreasing for $x>0$, and
$f(25/36)-1/32=55199/13436928>0$.
Thus either global sign of this complete rule fails to approximate $\tau$.
These statements do not cover signed differences of PSD matrices, nor
adding a freely chosen old approximant.

\subsection*{3. The projected PSD cone is a signed space}
Let $V_N=\operatorname{span}\{a_1,\ldots,a_N\}$, let $\Pi_N$ be its
orthogonal projection, and put $\phi_\ell=(I-\Pi_N)a_\ell$.
Define $T_N(\Gamma)=\sum_\ell q_\ell a_\ell$. Then, exactly,
\[
 \boxed{\{(I-\Pi_N)T_N(\Gamma):\Gamma\succeq0\}
 =\operatorname{span}\{\phi_\ell:\ell\in S_N,\ \ell>N\}.}
 \tag{96.9}
\]
One inclusion follows from support. For the other, choose $i,j\leq N$
with $[i,j]=\ell>N$; necessarily $i\ne j$. The two PSD rank-one matrices
$\Gamma_\pm=(e_i\pm e_j)(e_i\pm e_j)^T$ satisfy
\[
 T_N(\Gamma_\pm)=a_i+a_j\pm2a_\ell,
 \qquad (I-\Pi_N)T_N(\Gamma_\pm)=\pm2\phi_\ell.
 \tag{96.10}
\]
An arbitrary signed linear combination is therefore realized by adding
the corresponding matrices with nonnegative coefficients. This proof
covers unrestricted finite PSD mixtures; it does not assert the same
statement for a single rank-one square, fixed diagonal, normalization,
trace budget, or other additional constraints.

Equivalently, with an unrestricted old coefficient vector $d$,
$\sum_{n\leq N}d_n a_n+T_N(\Gamma)$ ranges over precisely
$\operatorname{span}\{a_\ell:\ell\in S_N\}$; the same holds for the
minus sign. The positivity is absorbed by old terms which are free to
cancel. This is an equality of approximation classes, not an error-decay
theorem. All old indices belong to $S_N$, but no prime $p>N$ belongs to it.
Also $S_N\subseteq\{ij:i,j\leq N\}$ since
$[i,j]=i(j/\gcd(i,j))$. No equality with NS95's divisor-completed space
is assumed.

\paragraph{Actual residual pairings.}
For any real $r\perp V_N$, the matrix
$M(r)_{ij}=\langle r,a_{[i,j]}\rangle$ has zero diagonal.
If it is nonzero, a nonzero off-diagonal entry $h$ gives the principal
minor $\left(\begin{smallmatrix}0&h\\h&0\end{smallmatrix}\right)$.
Thus $M(r)$ is indefinite and the two squares in (96.10) have pairings
$\pm2h$. If $M(r)=0$, all LCM observations vanish.
This statement applies in particular to the actual optimized residual,
but does not assert that any particular observation is nonzero.
It excludes an automatic one-sided numerator sign from PSD alone.

\paragraph{What remains after refitting.}
For complete coefficients $c=d-q$ with $d$ supported on $[1,N]$, let
$D(s)=\sum d_n n^{-s}$ and $z_s=-\zeta(s)>0$. The moment bound reads
\[
 |1-z_s(D(s)-Q(s))|\leq\rho_s,
 \quad \rho_s=s^2\|k_s\|\sqrt E,
 \quad \frac{1-\rho_s}{z_s}\leq D(s)-Q(s)
                           \leq\frac{1+\rho_s}{z_s}.
 \tag{96.11}
\]
For $E<(1-s)^2(2s-1)/s^6$, necessarily $D(s)>Q(s)$.
This is a necessary signed moment balance, not a sufficient estimate.
In particular (96.7) cannot be applied to refitted coefficients.

\subsection*{4. Stop condition and next input}
The LCM replacement repairs pointwise positivity, but as a complete rule
it cannot converge. The unrestricted old refit removes that obstruction
by making the projected PSD class exactly a signed space. Neither branch
supplies a cofinal arithmetic lower gain. No global Selberg-sieve no-go
or impossibility of a constrained, residual-dependent square is claimed.

The useful next input remains a joint estimate for the \emph{actual}
signed observations and their complete projected energy. One can keep
NS95's free signed class or use the full PR42 observation vector; a PSD
relabeling alone does not improve the estimate. A new restricted square
proposal must quantify its old-coefficient cancellation or trace budget,
prove its actual-target numerator and cost bounds together, and respect
NS83's logarithmic lower barrier. NS87 already controls the omitted tail
at the relevant scale; another finite efficiency table does not supply
the missing cofinal bound. G2, RH, and the original uniform Weil floor
remain open.
